% File: Volume-II-Applications/appendices/appI-numerical-evolution.tex

\chapter{Numerical Methods for Cosmological Evolution}
\label{app:cosmo-numerics}

This appendix summarizes numerical techniques relevant to solving the
cosmological evolution equations derived in Part~II of this volume.
We emphasize methods that are adapted to the lamphrodynamic field
equations and their interaction with entropy--vector dynamics.  Our
goal is to provide practical guidance for implementing numerical
solutions, rather than to develop a comprehensive theory of numerical
analysis.

\section{Discretization of Background Evolution}

We begin with the homogeneous and isotropic lamphrodynamic equations
derived in Chapter~4.  Let $a(t)$ denote the scale factor and
$\Psi(t)=(\PhiField(t),\vField(t),\SField(t))$ denote the spatially
homogeneous fields.  The coupled nonlinear ODE system takes the form
\[
  \frac{d}{dt}\Psi(t) = \mathcal{F}(\Psi(t),a(t))\,,
  \qquad
  \frac{d}{dt}a(t) = \mathcal{G}(\Psi(t),a(t))\,,
\]
where $\mathcal{F}$ and $\mathcal{G}$ are determined by the lamphrodynamic
action and its variational derivatives.

For time discretization, standard explicit and implicit Runge--Kutta
methods may be used.  Stability is improved by semi--implicit updates
for $\SField(t)$, reflecting entropy dissipation.  Adaptive step size
control is recommended due to the stiff coupling between $\vField$ and
$\SField$ during early epochs.

\section{Linear Perturbation Equations}

For linear perturbations, we decompose fields into Fourier modes
$\delta\Psi_k(t)$ satisfying
\[
  \frac{d}{dt}\delta\Psi_k(t)=\mathcal{L}_k\,\delta\Psi_k(t)\,,
\]
where $\mathcal{L}_k$ is the linearized operator.  The dominant numerical
issue is the scale dependence introduced by entropy--vector couplings,
leading to nontrivial dispersion.  Implicit schemes or operator
splitting techniques reduce numerical stiffness.

Spectral methods (Chebyshev, Fourier) are natural for periodic domains.
Finite--difference methods are viable on comoving grids, provided
appropriate boundary conditions and resolution criteria are enforced.

\section{Mode Coupling and Transfer Functions}

Transfer functions are computed by evolving perturbations from early
times to the present.  In \rsvpabbr{}, the lamphrodynamic contributions
modify both scalar and vector channels.  Numerical integration must
account for:
\begin{itemize}
  \item entropy--induced modification of gravitational potentials,
  \item cross--coupling terms in $\vField$ and $\SField$,
  \item altered growth rates compared to $\Lambda$CDM.
\end{itemize}
Parallelization across $k$--modes is straightforward and recommended.

\section{Weak Lensing and Line--of--Sight Integrals}

Weak lensing observables rely on line--of--sight integration of
perturbations along null geodesics.  In \rsvpabbr{}, null geodesics are
modified by entropy--vector contributions (Chapter~3).  Numerical
integration should incorporate:
\begin{itemize}
  \item modified geodesic equations,
  \item entropy--weighted optical depth,
  \item anisotropic stress contributions.
\end{itemize}
Adaptive sampling along the line of sight improves accuracy near sharp
features in $\SField(t)$ and $\vField(t)$.

\section{Boltzmann Solvers and Existing Software}

Standard Boltzmann solvers (CAMB, CLASS) can be extended to include
lamphrodynamic terms.  This requires:
\begin{itemize}
  \item modification of background equations,
  \item extension of perturbation modules to include $(\PhiField,\vField,\SField)$,
  \item additional source terms in line--of--sight integrators.
\end{itemize}
Validation against $\Lambda$CDM baselines is essential and provides a
useful diagnostic for implementation errors.

\section{Numerical Relativity and Nonlinear Regimes}

In strongly nonlinear regimes, numerical relativity methods become
necessary.  The lamphrodynamic field equations form a mixed system of
hyperbolic--parabolic type, where entropy diffusion plays a stabilizing
role.  Discontinuous Galerkin methods and finite--volume schemes with
shock capturing may be appropriate for studying entropic singularities.

\section{Parallelization and Performance}

Efficient parallelization exploits:
\begin{itemize}
  \item independent $k$--mode evolution,
  \item GPU acceleration for spectral transforms,
  \item domain decomposition for large--scale spatial grids.
\end{itemize}
Entropy diffusion reduces the stiffness of the system at late times,
but may increase computational cost during early evolution.

\section{Remarks}

The numerical implementation of \rsvpabbr{} remains an active area of
development.  Future work will combine Boltzmann solvers, numerical
relativity tools, and cosmological data pipelines to provide full
Bayesian inference for lamphrodynamic cosmology.  We return to these
questions in the concluding chapters of Volume~II.

